\chapter{Results}
\label{chap:results}

The sample means were close, with method B higher by
\SI{0.18}{\milli\metre} on average (\cref{tab:summary}).

\begin{table}[htbp]
  \centering
  \caption{Summary of simulated paired measurements.}
  \label{tab:summary}
  \begin{tabular}{lS[table-format=2.2]S[table-format=1.2]}
    \toprule
    Measure & {Mean} & {SD} \\
    \midrule
    Method A (\si{\milli\metre}) & 12.44 & 1.31 \\
    Method B (\si{\milli\metre}) & 12.62 & 1.35 \\
    Difference (\si{\milli\metre}) & 0.18 & 0.24 \\
    \bottomrule
  \end{tabular}
\end{table}

\begin{figure}[htbp]
  \centering
  \includegraphics[width=.72\textwidth]{method_comparison.pdf}
  \caption{Simulated measurements from methods A and B. The diagonal shows
  equality and is not a fitted model.}
  \label{fig:comparison}
\end{figure}

The figure and table are complementary: the figure shows individual pairs,
whereas the table provides compact numerical summaries.
